\section{Isomorphism with \texorpdfstring{$\Gamma_*$}{Gamma*}-Weighted Symmetric Functions}\label{sec:6}

This section is dedicated to proving that $\FHG$ is isomorphic to $\RGamma \otimes \LambdaG$ via Jucys--Murphy evaluation, analogously to Theorem \ref{thm:FH_lambda_iso}.

We introduce two filtrations. Let us say that $\sigma \in S_n$ \emph{moves} an element $m \in \{1,\ldots, n\}$ if $\sigma(m) \neq m$ (i.e. $m$ is not a fixed point of $\sigma$). The first filtration is obtained by placing an element $(\mathbf{g}, \sigma) \in \Gamma \wr S_n$ in degree $i$, where $i$ is the number of elements of $\{1, \ldots, n\}$ which are moved by $\sigma$. 
\begin{definition}
The \emph{moving filtration}, $\mathcal{F}_{mov}^i$, is the family of $\RGamma$-submodules of $\FHG$ spanned by $K_{\bm\hat{\bs \mu}}$ with $|\bs \mu| + l(\bs \mu) \leq i$.
\end{definition}

\begin{proposition} \label{moving_filt_prop}
The moving filtration is an algebra filtration. In the associated graded algebra, we have the following equation:
\[
K_{\bm\hat{\bs \mu}}K_{\bm\hat{\bs \nu}} = \left( \prod_{c \in \Gamma_*} \prod_{i \geq 1} \frac{m_i(\bs\lambda(c))!}{m_i(\bs\mu(c))!m_i(\bs\nu(c))!}\right) K_{\bm\hat{\bs \lambda}},
\]
where $m_i(\bs\lambda(c)) = m_i(\bs\mu(c)) + m_i(\bs\nu(c))$.
\end{proposition}

\begin{proof}
If $\sigma \in S_n$ has reduced cycle type $\mu$, the number of elements moved by $\sigma$ is $|\mu| + l(\mu)$. Similarly if $(\mathbf{g}, \sigma) \in \Gamma \wr S_n$ has fully-reduced cycle type $\bs\mu$, then the number of elements moved by $\sigma$ is $|\bs \mu| + l(\bs \mu)$. Note that this corresponds to the condition for membership in $\mathcal{F}_{mov}^i$, where $i = |\bs \mu| + l(\bs \mu)$. 

Fix $\bs \mu$ and $\bs \nu$ with $|\bs \mu| + l(\bs \mu) = i$ and $|\bs \nu| + l(\bs \nu) = j$. Consider $K_{\bm\hat{\bs \mu}}K_{\bm\hat{\bs \nu}}$. To compute this product, we pass to $Z(\mathbb{Z}\Gamma \wr S_n)$ for large $n$. A term in $X_{\bm\hat{\bs \mu}}$ is of the form $(\mathbf{g}, \sigma)$, where $\sigma$ moves $i$ elements. Similarly a term in $X_{\bm\hat{\bs \nu}}$ is of the form $(\mathbf{h}, \rho)$, where $\rho$ moves $j$ elements. Now, $(\mathbf{g}, \sigma) \cdot (\mathbf{h}, \rho) = (\mathbf{g}\sigma(\mathbf{h}), \sigma\rho)$. The permutation $\sigma \rho$ is obtained by moving $j$ elements, and then moving $i$ elements. This means that $\sigma \rho$ moves at most $i+j$ elements, with equality if and only if the elements moved by $\sigma$ and $\rho$ are disjoint. This proves the assertion that $\mathcal{F}_{mov}^i$ defines a filtration.

If $\sigma$ and $\rho$ move disjoint elements, then $\sigma(\mathbf{h}) = \mathbf{h}$ because the $m$-th factor in $\mathbf{h}$ is equal to $1$ if $m$ is not moved by $\rho$ (so in particular if $m$ is moved by $\sigma$). Hence the cycles (of size larger than 1) of $(\mathbf{g}, \sigma) \cdot (\mathbf{h}, \rho)$ are the constituent cycles of $(\mathbf{g}, \sigma)$ and $(\mathbf{h}, \rho)$. This shows that in the associated graded algebra, $K_{\bm\hat{\bs \mu}}K_{\bm\hat{\bs \nu}}$ is equal to a multiple of $K_{\bm\hat{\bs \lambda}}$. The multiple is equal to the number of ways to split the cycles of $\bs \lambda$ among $\bs \mu$ and $\bs \nu$.
\end{proof}


The second filtration is obtained by placing an element $(\mathbf{g}, \sigma) \in \Gamma \wr S_n$ in degree $i$, where $i$ is the smallest number of factors required to express $\sigma$ as a product of transpositions. 
\begin{definition}
The \emph{transposition filtration}, $\mathcal{F}_{tsp}^i$, is the family of $\RGamma$-submodules of $\FHG$ spanned by $K_{\bm\hat{\bs \mu}}$ with $|\bs \mu| \leq i$.
\end{definition}




\begin{proposition} \label{length_filtration_prop}
The transposition filtration is an algebra filtration. \end{proposition}

\begin{proof}
Similarly to the other filtration, we pass to $Z(\mathbb{Z}\Gamma \wr S_n)$ for $n$ sufficiently large. Fix $\bs \mu$ and $\bs \nu$ with $|\bs \mu| = i$ and $|\bs \nu| = j$. A term in $X_{\bm\hat{\bs \mu}}$ is of the form $(\mathbf{g}, \sigma)$, where $\sigma$ may be written as the product of $i$ transpositions. Similarly a term in $X_{\bm\hat{\bs \nu}}$ is of the form $(\mathbf{h}, \rho)$, where $\rho$ is a product of $j$ transpositions. Then, $(\mathbf{g}, \sigma) \cdot (\mathbf{h}, \rho) = (\mathbf{g}\sigma(\mathbf{h}), \sigma\rho)$. The permutation $\sigma \rho$ may be written as a product of $i+j$ transpositions (although possibly fewer). We conclude that $\mathcal{F}_{tsp}^i$ defines a filtration.
\end{proof}

\begin{remark}
Using \cite[Lemma 3.2]{FarahatHigman} and \cite[Proposition 2.9]{Wang1}, one can show that the structure constants of the $K_{\bm\hat{\bs\mu}}$ elements in the associated graded algebra are integers rather than arbitrary elements of $\RGamma$. We will not need this fact.
\end{remark}




\subsection{Jucys--Murphy Elements}
Analogously to the JM elements associated to symmetric groups, Pushkarev \cite{Pushkarev} and Wang \cite{Wang2} independently introduced the following elements of $\mathbb{Z} \Gamma \wr S_n$:
\[
L_j =
\sum_{i<j} \sum_{g \in \Gamma} g^{(i)} (g^{-1})^{(j)} (i, j).
\]
In the case where $\Gamma$ is the trivial group, we recover the JM elements for symmetric groups. However, these elements are not sufficient for our purposes, and we must introduce further generalisations.

\begin{definition}
For $c \in \Gamma_*$ and $j \leq n$, the \emph{wreath-product Jucys--Murphy elements} are
\[
L_j(c) = \sum_{i < j} \sum_{\substack{g_1, g_2 \in \Gamma \\ g_2g_1 \in c}} g_1^{(i)} g_2^{(j)} (i, j),
\]
which are elements of $\mathbb{Z} \Gamma \wr S_n$. We extend this notation by linearity in $c$: if $c = \sum_i  n_i c_i$ for some $n_i \in \mathbb{Z}$, we let
\[
L_j(c) = \sum_{i=1}^l n_i L_j(c_i).
\]
\end{definition}

In the case where $c = 1$, the condition $g_2g_1 \in c$ becomes $g_2 = g_1^{-1}$, so $L_j(1)$ recovers the definition of Pushkarev and Wang.
\begin{proposition}
Let $c \in \mathbb{Z}\Gamma_*$. We have
\[
L_j(c) = c^{(j)} L_j(1) = L_j(1) c^{(j)}.
\]
\end{proposition}
\begin{proof}
By linearity, it suffices to consider $c \in \Gamma_*$. To show the first equality, it suffices that for $i<j$,
\[
c^{(j)} \sum_{\substack{g_1, g_2 \in \Gamma \\ g_2g_1 = 1}} g_1^{(i)} g_2^{(j)} = \sum_{\substack{h_1, h_2 \in \Gamma \\ h_2h_1 \in c}} h_1^{(i)} h_2^{(j)},
\]
as we may right-multiply by $(i,j)$ and sum over $i$ less than $j$. We may set $g_2 = g_1^{-1}$ and use the fact that $c^{(j)}$ commutes with $g_1^{(i)}$ and $g_2^{(j)}$. The left hand side becomes
\[
\sum_{g_1 \in \Gamma} g_1^{(i)} (g_1^{-1})^{(j)} c^{(j)} = \sum_{g_1 \in \Gamma, k \in c} g_1^{(i)} (g_1^{-1} k)^{(j)},
\]
which agrees with the right hand side upon identifying $h_1 = g_1$ and $h_2 = g_1^{-1}k$. The second equality is similar, except that moving $c^{(j)}$ past the transposition $(i, j)$ turns it into $c^{(i)}$, which means that the manipulations take place in the $i$-index, rather than the $j$-index.
\end{proof}



\begin{definition}
For $c \in \Gamma_*$, let
\[
M_c^{(i,i+1)} = \sum_{\substack{g,h \in \Gamma \\ hg \in c}} g^{(i)}h^{(i+1)}.
\]
\end{definition}


The following lemmas have routine proofs, so we omit them.
\begin{lemma}
We have $M_c^{(i,i+1)} = c^{(i)} M_1^{(i,i+1)}  = c^{(i+1)} M_1^{(i,i+1)}$, and $M_c^{(i,i+1)}$ commutes with the transposition $s_i = (i, i+1)$.
\end{lemma}

\begin{lemma}[{\cite[Proposition 5]{Pushkarev}}]
We have that the $L_j(1)$ \textup{(}$j=1,\ldots,n$\textup{)} commute with each other. 
The $L_j(1)$ commute with with $\Gamma^n$. Furthermore
\[
s_i L_i(1) s_i + M_1^{(i,i+1)} s_i = L_{i+1}(1).
\]
\end{lemma}

The following lemma is easily proven by induction.
\begin{lemma} \label{iterative_swap_lemma}
We have the two identities
\[
s_i L_i(1)^r = L_{i+1}(1)^r s_i - \sum_{p=0}^{r-1} L_{i+1}(1)^p M_1^{(i,i+1)} L_i(1)^{r-1-p}
\]
and
\[
s_i L_{i+1}(1)^r = L_{i}(1)^r s_i + \sum_{p=0}^{r-1} L_{i}(1)^p M_1^{(i,i+1)} L_{i+1}(1)^{r-1-p}.
\]
\end{lemma}

\begin{proposition} \label{special_commuting_proposition}
Suppose $c_a, c_b \in \Gamma_*$. The following elements both commute with~$s_i$:
 \[
 L_i(1)^r c_a^{(i)}c_b^{(i+1)} + L_{i+1}(1)^r c_b^{(i)} c_a^{(i+1)}\text{ and }L_i(1)^r L_{i+1}(1)^r.
 \]
\end{proposition}

\begin{proof}
For the first part, we add the first equation in Lemma \ref{iterative_swap_lemma} right-multiplied by $c_a^{(i)}c_b^{(i+1)}$ to the second equation right-multiplied by $ c_b^{(i)}c_a^{(i+1)}$. Because $L_i(1)$ and $L_{i+1}(1)$ commute with $\Gamma^n$, they commute with $M_1^{(i,i+1)}$. So since $M_1^{(i,i+1)}  c_a^{(i)}c_b^{(i+1)} =  M_1^{(i,i+1)}c_b^{(i)}c_a^{(i+1)}$, the sums over the variable $p$ cancel out, leaving the required identity. For the second part it suffices to consider the case $r=1$ as the general case is recovered by raising $L_i(1)L_{i+1}(1)$ to a suitable power. Then applying Lemma \ref{iterative_swap_lemma} (specifically, the second equation, with $r=1$), we have
\begin{eqnarray*}
s_i L_i(1) L_{i+1}(1) &=& (L_{i+1}(1)s_i - M_1^{(i,i+1)}) L_{i+1}(1) \\
&=& L_{i+1}(1) (L_i(1)s_i + M_1^{(i,i+1)}) - M_1^{(i,i+1)} L_{i+1}(1) \\
&=& L_{i+1}(1)L_i(1)s_i.
\end{eqnarray*}
Since $L_i(1)$ and $L_{i+1}(1)$ commute, $s_i L_i(1)L_{i+1}(1) = L_{i}(1)L_{i+1}(1)s_i$ and we are done.\qedhere
\end{proof}


\begin{proposition} \label{gen_JM_elts_commute}
The wreath-product JM elements commute with each other: $L_{j_1}(c_1) L_{j_2}(c_2)=L_{j_2}(c_2) L_{j_1}(c_1)$ for any $c_1, c_2 \in \Gamma_*$ and $j_1, j_2 \in \{1,\ldots, n\}$.
\end{proposition}
\begin{proof}
We have that the elements $L_{j_1}(c_1) = L_{j_1}(1) c_1^{(j_1)}$ and $L_{j_2}(c_2) = L_{j_2}(1) c_2^{(j_2)}$ and $L_{j_1}(1), c_1^{(j_1)}, L_{j_2}(1), c_2^{(j_2)}$ commute pairwise.
\end{proof}






Because the wreath-product JM elements commute, it makes sense to evaluate a polynomial in them.
\begin{proposition}
There is a ring homomorphism
\[
ev_n^Q: Q^{\otimes n} \to \mathbb{Z} \Gamma \wr S_n,
\]
defined by 
\[
ev_n^Q(x_d^r(c)) = L_d(1)^r c^{(d)}.
\]
Moreover, the image of $\Lambda_n(\Gamma_*)$ \textup{(}the $S_n$-invariants of $Q^{\otimes n}$\textup{)} is contained in the centre of $\mathbb{Z} \Gamma \wr S_n$.
\end{proposition}
\begin{proof}
Since the wreath-product JM elements commute pairwise, it suffices to show that $ev_n^Q$ is well-defined on each tensor factor of $Q^{\otimes n}$. The $d$-th factor of $Q$ has a basis $x_d^r(c)$ for $r \geq 0$ and $c \in \Gamma_*$ (where we require $c=1$ if $r=0$). This basis obeys the multiplication rule mentioned at the start of Section \ref{LambdaG_section} (where we expand the multiplication in $\mathbb{Z}\Gamma_*$ in terms of the structure tensor $A_{i,j}^k$)
\[
x_d^r(c_i) x_d^s(c_j) = x_d^{r+s}(c_ic_j) = \sum_k A_{i,j}^k x_d^{r+s}(c_k).
\]
Correspondingly, using Proposition \ref{gen_JM_elts_commute},
\begin{eqnarray*}
ev_n^Q(x_d^r(c_i))ev_n^Q(x_d^s(c_j)) &=& L_d(1)^r c_i^{(d)} L_d(1)^s c_j^{(d)} \\
&=& L_d(1)^{r+s} (c_ic_j)^{(d)} \\
&=& L_d(1)^{r+s} \sum_k A_{i,j}^k c_k^{(d)} \\
&=& \sum_k A_{i,j}^k ev_n^Q(x_d^{r+s}(c_k)).
\end{eqnarray*}
We conclude that $ev_n^Q$ is well defined. 

Now suppose that $P(x_1, \ldots, x_n) \in (Q^{\otimes n})^{S_n}$. To show $ev_n^Q(P)$ is central, it suffices to show that it commutes with a generating set of $\Gamma \wr S_n$. In particular, we may take $\Gamma^n$ together with the adjacent transpositions $s_i = (i,i+1)$ for $i \in \{1,\ldots, n-1\}$. Since $L_j(1)$ commutes with $\Gamma^n$, the same is true for $L_j(c) = c^{(j)}L_j(1)$, from which it follows that the entire image of $ev_n^Q$ commutes with $\Gamma^n$. Now we show that $ev_n^Q(P)$ commutes with $s_i$. Since $P \in (Q^{\otimes n})^{S_n}$, $P$ is in particular an element of
\[
Q^{\otimes (i-1)} \otimes (Q \otimes Q)^{S_2} \otimes Q^{\otimes (n-i-1)},
\]
and $(Q \otimes Q)^{S_2}$ is spanned by elements of the form $x_i^r(c) x_{i+1}^r(c)$ and $x_i^r(c_a) x_{i+1}^s(c_b) + x_i^s(c_b) x_{i+1}^r(c_a)$. So it is enough that applying $ev_n$ to these elements gives something that commutes with $s_i$. In the first case,
\[
ev_n^Q(x_i^r(c) x_{i+1}^r(c)) = L_i(1)^r c^{(i)} L_{i+1}(1)^r c^{(i+1)},
\]
which commutes with $s_i$ because $L_i(1) L_{i+1}(1)$ and $c^{(i)}c^{(i+1)}$ do. For the second case, without loss of generality, we assume $s \leq r$. Now,
\scriptsize
\begin{eqnarray*}
ev_n^Q(x_i^r(c_a) x_{i+1}^s(c_b) + x_i^s(c_b) x_{i+1}^r(c_a)) &=& L_i(1)^r c_a^{(i)} L_{i+1}(1)^s c_b^{(i+1)} + L_i(1)^r c_b^{(i)} L_{i+1}(1)^s c_a^{(i+1)}\\
&=& L_i(1)^s L_{i+1}(1)^s \left( L_i(1)^{r-s}c_a^{(i)} c_b^{(i+1)} + L_{i+1}(1)^{r-s}c_b^{(i)} c_a^{(i+1)} \right)
\end{eqnarray*}
\normalsize
is a product of elements that commute with $s_i$ from Proposition \ref{special_commuting_proposition}.
\end{proof}

\begin{definition} \label{ev_n_formal_def}
Let $ev_n^\Lambda: \LambdaG \to Z(\mathbb{Z}\Gamma \wr S_n)$ be the composition of the canonical map $\LambdaG \to \Lambda_n(\Gamma_*) \subseteq Q^{\otimes n}$ with $ev_n^Q: Q^{\otimes n} \to Z(\mathbb{Z}\Gamma \wr S_n)$.
\end{definition}

Finally, we combine our evaluation maps into a single map.
\begin{definition}
For each $n\in \mathbb{Z}_{\geq 0}$ there is a ring homomorphism
\[
ev_n: \RGamma \otimes \LambdaG \to Z(\mathbb{Z}\Gamma \wr S_n)
\]
defined as 
\[
ev_n( a \otimes f) = ev_n^{\mathcal{R}}(a) ev_n^{\Lambda}(f),
\]
using the maps $ev_n^{\mathcal{R}}$ and $ev_n^{\Lambda}$ from Theorem \ref{r_gamma_hom_thm} and Definition \ref{ev_n_formal_def}, respectively.
\end{definition}

\begin{theorem} \label{wreath_map_exists}
There is a unique homomorphism $\Psi: \RGamma \otimes \LambdaG \to \FHG$ extending the map in Corollary \ref{coeff_compatibility_cor} such that the following diagram commutes.
\begin{equation*}
\begin{tikzcd}
\RGamma \otimes_{\mathbb{Z}} \LambdaG    \arrow[r, "\Psi"]    \arrow[dr, "ev_n"]    &    \FHG    \arrow[d, "\Phi_n"]\\
&    Z(\mathbb{Z}S_n)
\end{tikzcd}
\end{equation*}
\end{theorem}
\begin{proof}
Corollary \ref{coeff_compatibility_cor} constructs $\Psi$ on $\RGamma$, so it remains to construct $\Psi$ on $\LambdaG$. To do this, we pass to complex coefficients, and then show that the map is defined integrally. Proposition \ref{complex_bas_change_prop} shows
\[
\mathbb{C} \otimes \LambdaG = \mathbb{C} \otimes \Lambda^{\otimes |\Gamma^*|}
\]
by changing basis in $\mathbb{C}\Gamma_*$ from conjugacy-class sums to primitive central idempotents $e_\chi$. To compute $ev_n$ on the tensor factor of $\Lambda$ corresponding to $\chi \in \Gamma^*$, we must evaluate symmetric functions in the variables
\[
L_j(e_\chi) = L_j(1) e_\chi^{(j)}.
\]
Since $\Lambda$ is freely generated by the elementary symmetric functions $e_r$, it suffices to consider the evaluation of $e_r$ at the $L_j(e_\chi)$. Momentarily ignoring the $\Gamma^n$ parts in $L_j(e_\chi)$, we see that the transpositions involved in $L_j(e_\chi)$ are the same as in the JM elements $L_j$ for the symmetric group. So the $S_n$ parts are the same as in Theorem \ref{thm:elem_JM}. Lemma \ref{cycle_label_product} tells us that the labels of intersecting cycles multiply. But each label is the idempotent $e_\chi$. So we get a sum of elements whose reduced cycle type has size $r$, and all of whose cycles are labelled by $e_\chi$ (which is a linear combination of the usual labels in $\Gamma_*$). We could write this
\[
e_r(L_1(e_\chi), \ldots, L_n(e_\chi)) = \sum_{|\mu| = r} X_{\mu_\chi}^{irr},
\]
where the meaning of $X_{\mu_\chi}^{irr}$ is analogous to $m_{\bs\lambda_\chi}$ from Proposition \ref{complex_bas_change_prop}. Changing bases of $\mathbb{C}\Gamma_*$ back from $e_\chi$ to conjugacy-class sums changes the labels back to conjugacy classes. This is a linear operation that is independent of $n$. Hence we obtain some element
\[
e_r(L_1(e_\chi), \ldots, L_n(e_\chi)) = \sum_{|\bs\nu| = r} R_{r, \bs\nu}^\chi X_{\bs\hat{\bs \nu}},
\]
where $R_{r, \bs\nu}^\chi$ is independent of $n$, and the multipartition $\bs\nu$ is indexed by $\Gamma_*$ (the sizes of the cycles of the elements are unchanged by changing the type so we still sum over multipartitions of size $r$). Thus for $e_r \in \Lambda(\chi)$, we may take
\[
\Psi(e_r) = \sum_{|\bs\nu| = r} R_{r, \bs\nu}^\chi K_{\bm\hat{\bs \nu}},
\]
which guarantees that $ev_n = \Phi_n \circ \Psi$.

To see that the image of $\Psi$ on $\LambdaG$ is contained in $\FHG$ rather than $\mathbb{C}\otimes \FHG$, note that we can tell whether an element $f \in \mathbb{C} \otimes \FHG$ is contained in $\FHG$ by applying $\Phi_n$ for sufficiently large $n$ and checking whether the coefficients of the conjugacy class sums $X_{\bs\mu}$ are integers. This would imply that the coefficient of $K_{\bs\mu}$ in $f$, which is an element of $\mathbb{C} \otimes \mathcal{R}$, is integer-valued, i.e. an element of $\mathcal{R}$. On the other hand, we already know the image of $ev_n$ on $\LambdaG$ is contained in $Z(\mathbb{Z}\Gamma \wr S_n)$ and hence integral.

\end{proof}
